\documentclass[10pt,a4paper]{amsart}
\usepackage[margin=19mm]{geometry}
\usepackage{amsmath,amssymb,mathtools}
\usepackage[scr=rsfs,cal=euler]{mathalfa}
\usepackage{xfrac}
\usepackage[unicode,hidelinks,bookmarks=false]{hyperref}
\newcommand{\ie}{i.e.\ }
\newcommand{\cf}{cf.\ }
\newcommand{\resp}{respectively\ }
\newcommand{\Subsection}{\subsection}
\providecommand{\nobreakdash}{\nobreak\hbox{-}}
\setlength{\emergencystretch}{2em}
\multlinegap0pt
\usepackage{amssymb}
\usepackage{dsfont}
\newtheorem{proposition}{Proposition}
\title{A Ramanujanesque Suite}
\author{Zachary P. Bradshaw, C. Vignat}
\date{Source: arXiv:2605.08484v1 (CC BY 4.0)}
\begin{document}
\maketitle

\noindent\emph{Source and licence.} A Ramanujanesque Suite. Version arXiv:2605.08484v1 (CC BY 4.0), distributed by arXiv under the Creative Commons Attribution 4.0 International licence, http://creativecommons.org/licenses/by/4.0/, as recorded in the arXiv licence metadata for this exact version and shown on the abstract page https://arxiv.org/abs/2605.08484v1. Full licence notice: \texttt{licenses/arxiv-2605.08484-CC-BY-4.0.txt}.\par\smallskip

\noindent\textit{Editorial scope.} Counted unit: the article's proposition proposition (source0.tex lines 485--496). The article's complete original proof of it is delivered with it (source0.tex lines 497--532, one \texttt{proof} environment of 36 lines). No mathematics is written, repaired or re-derived: the statement and the proof are the article's own lines, in the article's own notation, and no retained line is substituted or re-flowed.  No further source-local context is needed: every cross-reference of every retained line resolves inside the two delivered spans.  Nothing was omitted from any retained span, so the package carries no omission record.  No retained line contains a citation, so the wrapper carries no bibliography and no bibliography entry.  No retained line contains a figure environment, an image-inclusion command or drawing code, so no image is delivered and the package declares no asset dependency.  Editorial formatting only, in addition: an AMS \texttt{amsart} preamble that restates the article's own declarations and macros used by the retained lines, namely the environments \texttt{proposition} on separate counters, together with the standard packages those lines need.  The retained environments print in this document as Proposition 1 in order of appearance, whereas the article prints the same items with the numbering of its own full document.  The complete native source of the article as distributed in the arXiv e-print for this exact version is bundled unchanged as \texttt{sources/200154/source0.tex}.
\begin{proposition}
For an arbitrary integer $q \ge 1,$ the following identity holds:
\begin{equation}
\log2 \sum_{l=1}^q
\sum_{k\ge2}\frac{\left(-1\right)^{k}}{k
\log\left(k2^l\right)}
+q\log^2(2)\sum_{k\ge2}\frac{1}{k
\log\left(2k\right)
\log\left(k2^{q+1}\right)}=\sum_{l=1}^q \frac{1}{l+1}.
\end{equation}

\end{proposition}

\begin{proof}
For $q=1,$ the partial fraction decomposition 
\begin{equation}
\nonumber
\frac{\log\left(2\right)}{\log\left(2k\right)\log\left(4k\right)}=\frac{1}{\log\left(2k\right)}-\frac{1}{\log\left(4k\right)}
\end{equation}
produces
\begin{equation}
\nonumber
\sum_{k\ge2}\frac{\left(-1\right)^{k}}{k\log\left(2k\right)}+\log2\sum_{k\ge2}\frac{1}{k\log\left(2k\right)\log\left(4k\right)}
=\sum_{k\ge2}\frac{\left(-1\right)^{k}}{k\log\left(2k\right)}+\sum_{k\ge2}\frac{1}{k\log\left(2k\right)}-
%\sum_{k\ge2}
\frac{1}{k\log\left(4k\right)}.
\end{equation}
Combining the first two series produces the series $\sum_{k\ge1}
\frac{1}{k\log\left(4k\right)}$ that telescopes with the third series to  produce $\frac{1}{2\log(2)}$ so that
\begin{equation}
\nonumber
\log2\sum_{k\ge2}\frac{\left(-1\right)^{k}}{k\log\left(2k\right)}+\log^{2}2\sum_{k\ge2}\frac{1}{k\log\left(2k\right)\log\left(4k\right)}=\frac{1}{2}.
\end{equation}
Similarly, in the case $q=2,$ the partial fraction decomposition 
\begin{equation}
\nonumber
\frac{\log\left(4\right)}{\log\left(2k\right)\log\left(8k\right)}=\frac{1}{\log\left(2k\right)}-\frac{1}{\log\left(8k\right)}
\end{equation}
produces
\begin{align}
\sum_{k\ge2}&\frac{\left(-1\right)^{k}}{k\log\left(2k\right)}+\sum_{k\ge2}\frac{\left(-1\right)^{k}}{k\log\left(4k\right)}+\log4\sum_{k\ge2}\frac{1}{k\log\left(2k\right)\log\left(8k\right)}\nonumber
\\
\nonumber
&=\sum_{k\ge2}\frac{\left(-1\right)^{k}}{k\log\left(2k\right)}+\sum_{k\ge2}\frac{\left(-1\right)^{k}}{k\log\left(4k\right)}+\sum_{k\ge2}\frac{1}{k\log\left(2k\right)}-\sum_{k\ge2}\frac{1}{k\log\left(8k\right)}.
\end{align}
The first and third series combine into $\sum_{k\ge1}\frac{1}{k\log\left(4k\right)}$ which in turn combines with the second series to produce $\frac{1}{\log 4}+\sum_{k\ge1}\frac{1}{k\log\left(8k\right)}.$ This series telescopes with the last one, leaving constant terms
$\frac{\log 2}{\log 4}+\frac{\log 2}{\log 8}=\frac{1}{2}+\frac{1}{3}$ as desired.
The proof of the general case follows the same pattern and is omitted.
\end{proof}

\end{document}
